\section{Discussion}
\label{sec:discussion}

This section frames what this paper establishes, situates the
deployment-safety theorem in the broader alignment-research
program, and enumerates the open questions deferred to follow-up
work.

\subsection{What this paper establishes}
\label{sec:disc_establishes}

Theorem~\ref{thm:deployment_safety} produces a structurally
defensible deployment-safety claim with three components:

\paragraph{An operationally checkable conditional theorem.}
Under the eleven invariants $I_1$--$I_{11}$, the canonical
tripartite substrate identification, the cooperative-overlap
regime, the causally-grounded inner-alignment condition, continuous
SPRT monitoring, and the SA1 surrogate-adequacy assumption, the
Goodhart slack between proxy and operational truth is bounded by a
constant intensive in the system's absolute capability magnitude.
Each conditional is named, ledger-observable (or, for SA1,
empirically testable), and operationally implementable via the
deployment-tooling specification of \S\ref{sec:deployment_tooling}.

\paragraph{Capability-magnitude independence.}
The bound's intensivity is the structural inversion of
capability-estimation-centered safety paradigms
(\S\ref{sec:wrong_frame}). Whatever capabilities a deployed system
actually possesses, the Goodhart slack stays bounded as long as
the operational invariants hold; the deployment claim scales past
any capability level the deployment's evaluation methods can
characterize. This is the punchline the title promises: a
\emph{provably safe regime for capability-unbounded deployment}.

\paragraph{Named residual structure.}
The four residuals (R1--R4 of Theorem~1) are explicit components of
the claim, not hidden caveats.  Operators can verify whether their
deployment falls outside each residual; the residuals collectively
name the boundaries of what the deployment claim covers.

The structural shape of the claim is a multi-conditional theorem
with each condition visible, each empirical assumption named, and
each scope restriction named.  This paper's contribution is not a
sweeping safety guarantee; it is the demonstration that a sweeping
guarantee is not necessary for operational deployment safety,
provided the conditions are structurally chosen to make the bound
capability-magnitude-independent.

\subsection{Connection to the alignment-research program}
\label{sec:disc_alignment}

This paper sits in the Goal-Frontier Maximization sequence
(\S\ref{sec:composition}) as the deployment-side closure of the
machinery developed across the prior GFM sequence. Three structural relationships
to the broader alignment-research program:

\paragraph{Inverting capability estimation.}
Standard alignment-safety paradigms place capability estimation at
the center: predict what the system can do, then bound the
consequences. This paper inverts this: name what the deployment can
operationally measure, prove that the bound holds under those
measurements regardless of underlying capability. The operational
invariants are upstream of capability estimation; the safety
guarantee follows from the regime, not from estimation accuracy.

This inversion is structurally available because the GFM machinery
(the prior GFM sequence) provides intensive bounds on the source-paper
quantities ($\rext$, $\rho_{\min}^{\mathrm{cross}}$, $\betaL$,
$\HHI$, $\epsnonres$). This paper's contribution is showing that
their composition is also intensive --- a non-trivial claim that
required substantive new content (the eleven invariants, the
Lemma~5 family, the cooperative-anchoring property).

\paragraph{The cooperative-anchoring stabilizing cascade.}
\S\ref{sec:cooperative_anchoring} establishes that under (C4),
optimization pressure on cooperative outputs is locally rational
toward preserving the substrate-exclusive verification layer. This
is \citep[Phase~Redundancy]{lasser2026phase}'s stabilizing cascade extended to the substrate level:
the verification infrastructure is not merely defended against
optimization pressure but is reinforced by it under proper
substrate identification.

The stabilizing cascade is conditional on (C4), basin entry, and
the eleven invariants. We do not claim alignment pressure is
universally reversed; we claim that the deployment dynamics admit
a defensible attractor where preservation is locally rational. The
mirror destabilizing cascade (\S\ref{sec:destabilizing_cascade})
is a real failure mode that operators must guard against.

\paragraph{Inner alignment as conditioning, not solution.}
This paper does not solve the inner-alignment problem. Condition
(C4) is conditioned, not discharged: the agent's effective
objective must realize causally-grounded cooperative-outcome
value, and vanilla RLHF does not deliver this. Achieving (C4)
requires training-time discipline (delayed outcome feedback,
adversarial fake-verification examples, process supervision tied
to real attestations). These are research directions in their own
right.

Operators using deployments that meet (C4) can invoke This paper's
theorem; operators using deployments that do not meet (C4) cannot.
The deployment claim does not produce safety where (C4) fails; it
produces safety where (C4) holds.

\subsection{Open questions deferred to companion work}
\label{sec:disc_companion}

Two structural items are deferred to a companion paper:

\paragraph{The channel-mediation conjecture.}
``All $\Delta r_K$ realizes through observable channels:
subsumption, redundancy loss, production concentration, trade-flow
concentration, yielding a structural bound $\Delta r_K \leq L_H
\cdot \HHI + L_\rho \cdot \Delta\rho + \ldots$.''  If proved, this
conjecture would close the channel-orthogonal residual (R1):
strategies that achieve $\Delta r_K \geq \rext$ without producing
any of the four monitored channel deviations would be impossible.
The current paper acknowledges R1 as a real residual; the
companion paper would establish whether R1 is empty or genuinely
contains undetectable strategies.

\paragraph{The HHI--$\Delta r_K$ structural relationship.}
The conjectural relationship between trade-flow concentration and
post-domination internal-rate efficiency. If proved, this would
either establish SA1 as a derived consequence (closing the
empirical-adequacy assumption) or refute it (forcing replacement
of $I_5$ with a different invariant). The empirical validation
pathway of \S\ref{sec:sa1_empirical} is the natural companion-paper
methodology.

The two open questions are related but distinct: the channel-
mediation conjecture closes R1; the HHI--$\Delta r_K$ relationship
closes the SA1 dependence. This paper's deployment claim is
strengthened by either; both would produce a deployment claim
conditional only on (C1)--(C5), not on SA1.

\subsection{Open questions deferred indefinitely}
\label{sec:disc_indefinite}

Four items are not This paper's scope and require their own research
programs:

\paragraph{The welfare-truth bridge.}
Stance S0 takes operational truth $\Ttruth = \volRwin$ as the
deployment claim's reference. The relationship between operational
truth and welfare-relevant truth ($\Ttruth_{\mathrm{welfare}}$, in
the sense of broader human flourishing) is not addressed. the Microfoundation paper~\citep{lasser2026micro}
\S goodhart's $T$-failure-modes remark establishes that the
operational and welfare-relevant truths can diverge; bridging them
is its own research program.

\paragraph{The monolithic-agent action-partition refinement.}
$I_4$ (subsumption frequency) and $I_5$ (HHI on trade flow)
presuppose a clean partition of the agent's action space. Monolithic
LLM-style agents without explicit subsumption operators may not
satisfy this. Extending the deployment claim to monolithic agents
requires either action-space-partition machinery (defining what
``subsumption'' and ``trade flow'' mean for a monolithic agent's
output stream) or alternative invariants that cover the same
structural roles.

\paragraph{Trusted-setup details for Pedersen commitments.}
This paper's verification infrastructure inherits \citep[Exogenous~Verification]{lasser2026exo}'s
Pedersen-commitment machinery, which requires a trusted setup.
The trusted-setup ceremony's security relies on at least one
honest participant; deployments where all participants are
compromised invalidate the commitment scheme. This is a real
operational risk that the deployment-tooling specification flags
but does not solve. Independent work on cryptographic primitives
that don't require trusted setup (or have weaker trust assumptions)
is the relevant research direction.

\paragraph{Empirical threshold calibration.}
The thresholds $(\theta_1, \ldots, \theta_{11})$ in the operational
invariants are policy choices. this paper commits to a methodology
(deployment-specific calibration via the deployment-tooling
specification) but does not fix specific numerical values. Real
deployments will need to choose thresholds based on the deployment's
specific threat model, capability profile, and operational
constraints. Empirical work on threshold calibration in real
deployments is a substantial follow-on.

\subsection{The structural significance of This paper's contribution}
\label{sec:disc_significance}

We close with three observations on what this paper contributes to
the alignment-research program.

\paragraph{Existence of a provably safe regime.}
Theorem~1 establishes that the regime is non-empty: under the
named conditions, there exist deployments for which the
deployment-safety bound binds. Scenario~1 of
\S\ref{sec:scenario_clean} shows what such a deployment looks like
operationally. The conditions are non-trivial, but they are
operationally achievable.

\paragraph{Capability-magnitude independence as the structural
inversion.}
The intensivity property of the bound is the structural inversion
of capability-estimation-centered safety. Once capability
estimation is no longer load-bearing, the deployment-safety
guarantee scales past any capability level. The inversion is
available because the GFM machinery's intensive bounds compose to
intensive bounds; without the GFM machinery, the inversion does not
work.

\paragraph{The conditional theorem as the appropriate structural
shape.}
The argument that conditional theorems with named conditions are
\emph{less} valuable than unconditional theorems mistakes the
nature of safety claims. Unconditional safety claims under
realistic assumptions about capability-unbounded systems are not
available; conditional safety claims with operationally checkable
conditions are. This paper's contribution is to demonstrate that the
conditional-theorem structure produces a defensible deployment
claim, not just an academic exercise.

The deployment-safety guarantee is conditional, but each condition
is named, ledger-observable, and operationally implementable. The
residuals are explicit. The empirical assumptions are testable.
This is the structural shape we believe deployment-safety claims
must take in the era of capability-unbounded systems --- not because
weaker claims are aesthetically preferable, but because they are
the strongest defensible claims the structural apparatus supports.
